\section{附录}

本文完整源程序、输入数据和中间结果随支撑材料提交。本附录给出关键
求解过程的伪代码，用于说明正文模型与程序实现之间的对应关系。伪代码
只保留主要决策逻辑，不展开数据读取、单位换算和结果写出等工程细节。

\subsection{问题二离散开停机调度}

问题二在给定场景和目标日产量下，只需确定开机小时集合。由于题设未考虑
启停成本、爬坡约束和跨时段库存，日内各小时除开机小时数约束外相互独立，
因此可采用净成本增量排序法求解。

\begin{enumerate}[label=\textbf{步骤~\arabic*.},leftmargin=4em]
  \item 输入风电、光伏、常规负荷和分时电价序列，给定目标日产量 $q$。
  \item 令所需开机小时数 $H=q/3$。
  \item 对每个时段 $t$，分别计算停机状态 $u_t=0$ 与满开状态 $u_t=1$
        下的园区用电功率。
  \item 根据功率缺额计算两种状态下的购电功率和上网功率：
        \[
          P_t^{\mathrm{buy}}=\max(P_t^{\mathrm{use}}-P_t^{\mathrm{re}},0),
          \quad
          P_t^{\mathrm{sell}}=\max(P_t^{\mathrm{re}}-P_t^{\mathrm{use}},0).
        \]
  \item 分别计算两种状态下的风光度电成本、设备运维成本、分时购电成本
        和上网收入抵扣。
  \item 计算该时段由停机转为开机的净成本增量
        $\Delta C_t=C_t(u_t=1)-C_t(u_t=0)$。
  \item 将 24 个时段按 $\Delta C_t$ 从小到大排序，选择前 $H$ 个时段开机，
        其余时段停机。
  \item 根据所得 $u_t$ 序列重新计算逐时功率平衡、三项绿电直连指标和吨氨成本。
\end{enumerate}

\subsection{问题三连续负荷率优化}

问题三将离散开停机变量 $u_t$ 替换为连续负荷率 $\alpha_t$，并设置
$0.1\leq\alpha_t\leq1$。程序中按日场景建立线性约束，并通过购售电互斥
条件避免同一时段同时买电和卖电。

\begin{enumerate}[label=\textbf{步骤~\arabic*.},leftmargin=4em]
  \item 对每个风光组合场景 $s$ 和日产量 $q$，建立变量
        $\alpha_t$、$P_t^{\mathrm{buy}}$、$P_t^{\mathrm{sell}}$ 和互斥变量 $z_t$。
  \item 添加负荷率约束：
        \[
          0.1\leq\alpha_t\leq1,\qquad t=1,2,\dots,24.
        \]
  \item 添加日产量约束：
        \[
          \sum_{t=1}^{24}3\alpha_t=q.
        \]
  \item 由 $P_t^{\mathrm{use}}=P_t^{\mathrm{load}}+41.5\alpha_t$
        建立逐时功率平衡：
        \[
          P_t^{\mathrm{re}}+P_t^{\mathrm{buy}}
          =P_t^{\mathrm{use}}+P_t^{\mathrm{sell}}.
        \]
  \item 通过大 $M$ 约束限制 $P_t^{\mathrm{buy}}$ 与
        $P_t^{\mathrm{sell}}$ 不能同时为正：
        \[
          P_t^{\mathrm{buy}}\leq Mz_t,\qquad
          P_t^{\mathrm{sell}}\leq M(1-z_t).
        \]
  \item 以日运行净成本最小为目标求解，得到连续负荷率序列。
  \item 汇总 24 场景和五个日产量档位的逐时结果，并折算全年成本和绿电指标。
\end{enumerate}

\subsection{问题四离网储能配置}

问题四分为无储能离网调度、能源自治装机估算和储能候选配置比较三部分。
其中储能配置不直接追求最大容量，而是在制氨产量、弃电量和吨氨成本之间
比较经济折中方案。

\begin{enumerate}[label=\textbf{步骤~\arabic*.},leftmargin=4em]
  \item 在无储能条件下，先以风光出力满足常规负荷。
  \item 若剩余功率不低于 $0.1P^{\max}$，则按式~\eqref{eq:p4-no-storage-alpha}
        确定制氨负荷率；否则制氨装置停机。
  \item 对所有场景逐时计算制氨量、弃电量和常规负荷缺额，并识别最大弃电场景。
  \item 根据全时段常规负荷保障条件，计算等比例扩容倍数 $k_{\min}$；
        同时在给定步长下搜索最小新增装机组合。
  \item 对最大弃电场景枚举候选储能功率和容量，建立充电功率、放电功率、
        荷电状态、制氨负荷率和弃电功率之间的日内平衡。
  \item 对每个候选配置求解离网调度，记录日产氨量、弃电量、常规负荷缺额、
        储能年化成本和吨氨成本。
  \item 在候选结果中选择经济折中配置，并将该配置用于 24 场景全年折算。
\end{enumerate}
